\documentclass{report}
\usepackage{amsmath,amssymb}
\setlength{\parindent}{0mm}
\setlength{\parskip}{1em}
\begin{document}
\begin{center}
\rule{6in}{1pt} \
{\large David Appelhans \\
{\bf Nested Iteration Range Decomposition: A scalable, low communication PDE solution method with numerical results on advection-diffusion problem.}}

99 Grant Ave #1 \\ Medford \\ MA 02155
\\
{\tt dappelh@us.ibm.com}\\
John Ruge\\
Tom Manteuffel\\
Steve McCormick\end{center}

Combining nested iteration, adaptive mesh refinement, and range
decomposition creates a powerful PDE solution method with reduced
communication costs compared to traditional solution methods. Range
decomposition partitions a PDE according to the right side of the
equation, rather than the domain. If the right side is approximately
equally distributed, for example from a coarse adaptive least squares
solve, then range decomposition equally distributes the error and is a
naturally load balanced algorithm. This algorithm is naturally suited for
large distributed memory systems where intra-node communication is
expensive but on node computational power is plentiful. The method will
be explained, a performance model will show the scaling benefits, and
numerical results for an advection-diffusion type PDE will be shown.


\end{document}
